% Generated draft: requires mathematical and source review.
\documentclass[11pt]{article}
\usepackage[margin=25mm]{geometry}
\usepackage{fontspec}
\setmainfont{Latin Modern Roman}
\usepackage{amsmath,amssymb,mathtools,unicode-math}
\setmathfont{Latin Modern Math}
\usepackage{enumitem,needspace,xurl,hyperref}
\setlength{\parindent}{0pt}
\setlength{\parskip}{5pt}
\setlength{\emergencystretch}{3em}
\newcommand{\N}{\mathbb{N}}
\newcommand{\Z}{\mathbb{Z}}
\newcommand{\Q}{\mathbb{Q}}
\newcommand{\R}{\mathbb{R}}
\newcommand{\C}{\mathbb{C}}
\newcommand{\F}{\mathbb{F}}
\newcommand{\A}{\mathbb{A}}
\newcommand{\PP}{\mathbb{P}}
\newcommand{\E}{\mathbb{E}}
\newcommand{\eps}{\varepsilon}
\newcommand{\dd}{\,\mathrm{d}}
\begin{document}
\section*{UnsolvedMath ID: 800282}
\textbf{Source problem code:} catalogue entry
\subsection{Definitions and mathematical statement}
Let \(M\) and \(N\) be strongly pseudoconvex CR manifolds of real dimension \(3\). Write \(H_M\subset TM\) and \(H_N\subset TN\) for their real rank-two CR distributions, with complex structures \(J_M\) and \(J_N\). Strong pseudoconvexity makes these distributions contact distributions. Choose smooth positive-definite inner products on the distributions compatible with their complex structures. The associated Carnot--Carath\'eodory distances are defined by
\[
d_M(p,q)=\inf_{\gamma}\int_0^1\|\dot\gamma(t)\|_{H_M}\,\dd t,
\]
where the infimum is over absolutely continuous curves from \(p\) to \(q\) satisfying \(\dot\gamma(t)\in H_M\) almost everywhere; define \(d_N\) analogously.

The source calls a homeomorphism \(f:M\to N\) conformal when it maps infinitesimal Carnot--Carath\'eodory spheres to infinitesimal spheres. One possible precise metric interpretation of this condition is the following. For each \(x\in M\), set
\[
L_f(x,r)=\sup_{d_M(x,y)=r}d_N(f(x),f(y)),\qquad
\ell_f(x,r)=\inf_{d_M(x,y)=r}d_N(f(x),f(y)),
\]
and require
\[
\limsup_{r\downarrow0}\frac{L_f(x,r)}{\ell_f(x,r)}=1
\qquad\text{for every }x\in M.
\]
This is an explicit candidate interpretation, not a quantified definition supplied by the source.

For a \(C^1\) map, the CR condition is
\[
\dd f_x(H_{M,x})\subseteq H_{N,f(x)},\qquad
\dd f_x(J_Mv)=J_N\dd f_x(v)
\quad(x\in M,\ v\in H_{M,x}).
\]
\textbf{Question (Puqi Tang; source Problem 82).} Is every conformal homeomorphism between strongly pseudoconvex CR \(3\)-manifolds CR?

\textbf{Unresolved qualifications.} The source does not quantify its infinitesimal-sphere condition or specify the regularity intended in the conclusion. It also imposes no orientation or CR-orientation condition and does not state whether anti-CR maps are to be excluded or included in the conclusion. These issues must be resolved before the candidate metric interpretation above can be treated as a definitive formulation. The source's broader introductory assertion that a Carnot--Carath\'eodory metric can always be defined on a CR \(3\)-manifold also requires qualification: the horizontal length construction needs an accessibility hypothesis to give finite distances between relevant points; strong pseudoconvexity supplies the local contact hypothesis in the actual question.
\subsection{Short English statement}
Puqi Tang asks whether a homeomorphism between strongly pseudoconvex CR manifolds of real dimension \(3\) that preserves infinitesimal Carnot--Carath\'eodory spheres must preserve the CR structures. The meaning of infinitesimal sphere preservation, the required regularity, and the treatment of anti-CR maps remain unspecified in the supplied source.
\subsection{Sources}
\begin{itemize}
\item Supplied local UnsolvedMath registry record, permanent ID \(800282\), exact source code \texttt{catalogue entry}, titled \emph{Several Complex Variables Problem List --- Problem 82}. The registry labels the problem open; this supplied label is not an independently verified mathematical or literature assessment.
\item Buzzard, Fornaess, Krantz, \emph{The Several Complex Variables Problem List} (1995), Problem 82. Supplied source URL: \url{https://arxiv.org/abs/math/9509202}. The supplied provenance records access on 2026-10-16 and extraction by \texttt{arxiv\_latex}; neither access nor source contents have been independently verified here. The extracted entry bears the label ``(New)'' and describes the question as a Liouville type assertion on CR \(3\)-manifolds.
\item The registry's proposer field names Buzzard, Fornaess, Krantz, whereas the extracted problem itself is signed Puqi Tang, with supplied contact \texttt{tang@math.purdue.edu}. The question is therefore attributed here to Tang within that collected problem list.
\end{itemize}
\end{document}
